\scanpage{46}
\mspage{133}

On a alors
\[
\begin{aligned}
H^1(X,F)&\simeq H^0(S,R^1\deleted{\unsure{f}}\unsure{p|f}_*(F))\\
H^1(X',F')&\simeq H^0(S,R^1p'_*(F'))
\end{aligned}
\]
et on voit encore que
\[
R^1p_*(F)\unsure{\text{---}|\to}R^1\deleted{\unsure{f}}p'_*(F')
\]
est iso (car \unsure{ils|il} se \unsure{calculent|calcule} fibre par fibre,
et \deleted{\unsure{alors ils}} on utilise l'hyp. sur les $f_s$),
donc induit un iso
\[
H^1(X,F)\text{---}H^1(X',F').
\]

$3^\circ$) Si $F$ est \deleted{\unsure{un}} abélien, on doit prouver
\[
H^i(X,F)\xrightarrow{\sim}H^i(X',F')
\]
pour tout $i$, et cela résulte encore de la
suite spectrale de Leray pour $(p,F)$ et
pour $(p',F')$.

b) Inversement, supposons que $f$ soit une
Hot-équivalence, et prouvons que les $f_s:X'_s\to X_s$
sont des Hot-équivalences. Je vais utiliser le

\uline{Lemme} Soit $S\in\mathrm{Ob}\,\mathrm{Cat}$, \inserted{$S$ \inserted{0-}connexe.} Alors il existe un
\begin{tabular}{|p{.90\linewidth}}
$\bar S$ \deleted{\unsure{asphé}} $\in\mathrm{Ob}\,\mathrm{Cat}$ \deleted{et}, $\bar S$ Hot-asphérique,
et $\bar S\longrightarrow S$ une Hot-fibration,
\end{tabular}

Voyons comment ceci implique le
th. 2 (partie réciproque) -- \unsure{OPS|O.P.S.} $S$ 0-connexe,
et \deleted{\unsure{on}} fait le chgt de base par $\bar S\to S$

% archive sheet number: 46
